\chapter{Physics and tooling}

The entries here are about the correctness apparatus rather than the physics
kernels. Two of them are cases where the test was wrong, which is worth recording precisely because a wrong test is more dangerous than
a wrong implementation: it either blocks correct work or, far worse, passes
when it should not.

\entry{The hash collision that would have counted neighbours twice}

\field{The pathology.} Spatial hashing is many to one by construction. Two
distinct grid cells can map to the same bucket. The neighbour search visits the
twenty seven cells around a particle and scans each one's bucket, so if two of
those twenty seven cells collide in the table, the particles in that bucket are
scanned twice and counted twice.

The reason this is nasty is that the usual defence does not work. Filtering
candidates by distance removes particles from unrelated cells, because they are
genuinely far away. It does nothing about double counting, because those
particles really are within the smoothing radius; they are simply being added
to the sum more than once. The density would come out high by an amount that
depends on the hash function rather than on the fluid, and the fluid would
still look like a fluid.

\field{Fix and why.} Each particle stores an exact cell identity packed into 64
bits as three 21 bit fields, and the bucket scan compares it before accepting a
candidate. This rejects both classes at once, at the cost of one 64 bit compare
per candidate. The alternative of deduplicating the twenty seven hashed buckets
before scanning would cost hundreds of comparisons per particle rather than
one per candidate.

\field{Verification.} This is the one defect in the project that was designed
against rather than discovered, so the verification carries the whole weight.
The density field produced by the hashed search is compared against a brute
force $O(N^2)$ host sum over identical positions, and the worst relative
difference across four thousand particles is \NeighbourSearchError{}, about what single
precision rounding gives. A single dropped or double counted neighbour
would move a density by several percent. A separate health check counts
particles sharing a bucket with a foreign cell and confirms the collision rate
is real but small, so the test is exercising the defence rather than a table so
sparse that no collisions occur.

\entry{A diagnostic that would have perturbed the run it measured}

\field{The problem.} The fused variant carries the leapfrog half step velocity,
which is not the velocity belonging to the current positions. The energy
diagnostic needs the synchronised one. The obvious implementation is to shift
the velocity array forward by half a step, take the measurement, and shift it
back.

\field{Why that is wrong.} In single precision, adding a quantity and then
subtracting it does not return the original bits. A run that happened to take a
diagnostic would therefore follow a very slightly different trajectory from one
that did not, which breaks the bit identity guarantee the resume path claims
and makes the energy curve partly a measurement of its own observer.

\field{Fix.} The half step correction is applied inside the reduction kernel,
which reads the velocity and the stored acceleration and forms $v + a\,\Delta
t/2$ on the fly without writing anything back.

\field{Verification.} A test runs eighty steps taking a full diagnostic after
every step, and asserts the final state is bit for bit identical to the same
run with no diagnostics at all.

\entry{The density gate that passed by measuring the wrong particles}

\field{Symptom.} None, which is the point. \texttt{v1.0.0} published an interior density error of
4.20 percent against a 5 percent limit, with an explanation about smoothing on a jittered lattice.
It was found by review for this release rather than by any failure.

\field{Root cause.} Two defects, one hiding the other. The gate chose interior particles by position
in lattice order and read densities in cell sorted order, so positions and densities belonged to
different particles, and the figures were close to the statistics of the whole block. With the index
fixed, the \texttt{v1.0.0} mass gave 7.61 percent, which fails that limit. And that mass was wrong: it
assumed a spacing the lattice does not have, and ignored the kernel's own lattice sum error.

\field{Consequence.} The equation of state turns excess density into pressure, so the
\texttt{v1.0.0} dam break interior started near 76 kPa against a hydrostatic 3.9 kPa. Maximum speed
after 100 steps was 3.695 m/s. With the calibrated mass it is 0.625 m/s.

\field{Fix.} The index is fixed, the mass comes from a kernel sum over the lattice laid out, and new
limits were committed before the code measuring them existed: the mean error is
\SphDensityErrorPercent{} percent, the RMS deviation \SphDensityRmsPercent{} percent and the maximum
deviation \SphDensityMaxPercent{} percent.

\field{Why it is in this report.} The gate read well, passed, and came with a plausible story. A
check that was never seen to fail for the right reason should not be trusted just because it passes.

\entry{A test that failed for a reason that was entirely my fault}

\field{Symptom.} The SPH kernel unit test
\texttt{KernelsAreContinuousAtTheSupportBoundary} failed on first run, while
every other kernel test passed.

\field{Root cause.} The test asserted that each kernel evaluated just inside its
support boundary was within $10^{-3}$ of zero, in absolute terms. These kernels
carry a $1/h^6$ normalisation: at a smoothing radius of $0.1$ the spiky
gradient peaks around $1.4 \times 10^5$ and the viscosity Laplacian around
$1.4 \times 10^6$. An absolute tolerance of $10^{-3}$ against quantities of that
magnitude is not a statement about continuity, it is a statement about nothing.
The implementation was correct throughout.

\field{Fix.} The test now compares each kernel at the boundary against its own
peak value, which is what continuity actually means for a quantity with a
scale, and separately asserts the kernels are exactly zero outside the support
so the transition is genuinely continuous rather than merely small.

\field{Lesson.} A tolerance without a scale is not a tolerance. Every numerical
gate in this project is now expressed relative to something: the magnitude of
the quantity, the machine epsilon, or a bound derived from the arithmetic.

\entry{A benchmark that recorded a random seed as a tile width}

\field{Symptom.} The tile width figure had four bars, one of which sat at a
tile width of \texttt{20260802}. That number is the random seed.

\field{Root cause.} The tile sweep opened the output file twice. It created a
truncating output stream to write the header, and then called the row appending
helper, which opens the same path independently in append mode. The header
stream stayed open and buffered for the duration of the sweep and flushed at
offset zero when it was finally destroyed, overwriting the beginning of the
file with the header text. What survived was a file whose first rows were
partly overwritten, so the column boundaries shifted and fields were read from
the wrong columns.

\field{Fix.} The header is written through the same append path as the rows,
with the file removed first, so only one handle to the file exists at any time.

\field{Why it is in this report.} Nothing failed. The benchmark ran, the CSV
parsed, the plot rendered, and the number was wrong. It was caught by looking
at the figure and asking what a tile width of twenty million meant. The general
form of this defect is that two independent handles to one file will corrupt
it, and buffering makes the corruption appear far from its cause.

\entry{LaTeX macro names cannot contain digits}

\field{Symptom.} The report failed to compile with
\texttt{Missing \textbackslash begin\{document\}} pointing at a line in the
generated macro file that looked entirely reasonable.

\field{Root cause.} A LaTeX control sequence is made of letters only. The
offending definition was

\begin{lstlisting}
\newcommand{\Poly6NormalisationError}{...}
\end{lstlisting}

\noindent
in which the name ends at the digit, so this defined a macro called
\texttt{Poly} followed by the literal text \texttt{6NormalisationError}. Literal
text in the preamble is what produces that error message, and it is reported at
the point of the text rather than at the mistake.

\field{Fix.} The three offending macros were renamed to spell their digits, and
the generator now rejects any macro name that is not purely alphabetic, so this
class of mistake fails at generation time with a message that names the
problem rather than at compile time with one that does not.

\entry{BibTeX put en dashes back after they had been removed}

\field{Symptom.} The dash gate failed on the compiled PDF, reporting seven en
dashes in the bibliography, after the source bibliography had been cleaned of
double hyphens and the gate had reported the sources clean.

\field{Root cause.} The bibliography style is what dashifies page ranges. The
\texttt{plain} style passes page fields through a function that converts any
run of hyphens into a TeX en dash ligature when it writes the intermediate
bibliography file. Cleaning the source had no effect because the source was
never the problem.

\field{Fix.} Page ranges are written with the word \texttt{to}, which the style
passes through untouched. This also makes the bibliography consistent with the
prose, which is required to write ranges the same way.

\field{Why this justifies the whole design of the gate.} The rule forbids a
character in the finished document, so the gate has to check the finished
document. A source only check would have reported clean while every reference
in the printed bibliography broke the rule. Extracting the text of the compiled
PDF and scanning that is what caught it, and it is the reason the gate runs
last in the report build rather than first.
